% =====================================================================
% Capítulo: Construcción del Sistema
% Fuentes: docs/tfg/{PLAN,PROGRESO,BITACORA,DECISIONES,TRAZABILIDAD}.md,
% MAPA-REFERENCIAS.md (solo como guía; no se nombran las referencias),
% AGENTS.md, .claude/{settings.json,hooks,skills,agents},
% apps/web/supabase/{schemas,migrations}, packages/cms/*,
% packages/function-middleware, apps/web/src. Estado: F3c cerrada.
% =====================================================================

% Estilo de los extractos de código de este capítulo: letra pequeña, marco
% y sustitución de los caracteres acentuados (listings no admite UTF-8
% multibyte de forma nativa).
\lstdefinestyle{pymekit}{
  basicstyle=\ttfamily\scriptsize,
  breaklines=true,
  frame=single,
  columns=fullflexible,
  keepspaces=true,
  showstringspaces=false,
  numbers=left,
  numberstyle=\tiny,
  numbersep=6pt,
  xleftmargin=1.2em,
  literate={á}{{\'a}}1 {é}{{\'e}}1 {í}{{\'i}}1 {ó}{{\'o}}1 {ú}{{\'u}}1
           {Á}{{\'A}}1 {É}{{\'E}}1 {Í}{{\'I}}1 {Ó}{{\'O}}1 {Ú}{{\'U}}1
           {ñ}{{\~n}}1 {Ñ}{{\~N}}1 {ü}{{\"u}}1 {¿}{{?`}}1 {¡}{{!`}}1
           {«}{{\guillemotleft}}1 {»}{{\guillemotright}}1
}

Este capítulo describe cómo se ha construido PymeKit: el entorno y las
herramientas, la estructura del código, qué parte se ha reutilizado y qué
parte es propia, la secuencia de incrementos (fases F0 a F3c), una selección
de extractos de código representativos, el uso de agentes de IA en el
desarrollo y los problemas más relevantes que se resolvieron.

% ---------------------------------------------------------------------
\section{Entorno y herramientas}
\label{sec:imp-entorno}

El desarrollo se ha realizado en un servidor Linux remoto, accedido por SSH,
con las herramientas de la tabla~\ref{tab:imp-herramientas}. La base de datos
local (PostgreSQL, autenticación y almacenamiento) se ejecuta en contenedores
con la CLI de Supabase; como el servidor tiene IP pública, sus puertos se
filtran explícitamente y la aplicación se consulta a través de un túnel SSH
(ADR-010).

\begin{table}[htbp]
  \centering
  \footnotesize
  \caption{Herramientas del entorno de desarrollo.}
  \label{tab:imp-herramientas}
  \begin{tabularx}{\textwidth}{@{}>{\raggedright\arraybackslash}p{3.4cm}X@{}}
    \toprule
    \textbf{Ámbito} & \textbf{Herramienta} \\
    \midrule
    Lenguaje y ejecución & TypeScript en modo estricto, Node.js (versión LTS) \\
    Monorepo & pnpm (con catálogo de versiones compartido) y Turborepo~\cite{turborepo} \\
    Aplicación & TanStack Start, Router y Query; React 19; Vite; Tailwind CSS 4 \\
    CMS & Hono, Drizzle ORM, Zod \\
    Base de datos & Supabase CLI (PostgreSQL, Auth, Storage) en contenedores Docker \\
    Pagos & Stripe en modo de pruebas y su CLI para reenviar \textit{webhooks} \\
    Pruebas & Vitest (unitarias), pgTAP~\cite{pgtap} (base de datos), Playwright~\cite{playwright} (E2E) \\
    Calidad & ESLint, Prettier, comprobación de tipos y de dependencias del monorepo \\
    Control de versiones & Git, GitHub (repositorio privado) y GitHub Actions; mensajes según Conventional Commits~\cite{conventionalcommits} \\
    Asistencia de IA & Claude Code~\cite{claudecode} con un \textit{harness} propio (sección~\ref{sec:imp-agentes}) \\
    Memoria & \LaTeX{} con la plantilla oficial de la escuela \\
    \bottomrule
  \end{tabularx}
\end{table}

La integración continua (\texttt{.github/workflows/ci.yml}) tiene tres
trabajos: \emph{Controles del TFG} (control de marca, bloqueante, e informe de
comentarios en inglés), \emph{TypeScript} (tipos, \textit{lint}, formato y
pruebas unitarias en cada \textit{push} a la rama principal y en cada \textit{pull request}) y
\emph{Supabase + E2E} (base de datos local, pgTAP, \textit{build} de pruebas y
Playwright), que por su coste solo se ejecuta en las \textit{pull requests}
hacia la rama principal o a mano.

% ---------------------------------------------------------------------
\section{Estructura del código}
\label{sec:imp-estructura}

La figura~\ref{fig:imp-arbol} muestra la estructura del repositorio. La
aplicación web es la única aplicación desplegable; todo lo demás son paquetes
que la componen o herramientas de verificación.

\begin{figure}[htbp]
  \centering
\begin{lstlisting}[style=pymekit, numbers=none, frame=none]
apps/
  web/                      TanStack Start + React 19 (una sola aplicación)
    src/routes/             rutas por fichero
      _marketing/           web pública: inicio, precios, blog, legales
      _authenticated/       área privada (sesión viva y MFA)
      admin/                consola: plataforma (super-admin) y cms/ (CMS)
      api/                  rutas de servidor: cms/$, billing/webhook, db/...
    src/lib/                server functions (*.functions.ts), servicios (*.server.ts)
    src/config/             configuración reutilizable (*.config.ts)
    supabase/               schemas/ (00-55), migrations/, tests/database (pgTAP)
  e2e/                      pruebas Playwright de la web y del CMS
packages/
  ui, i18n, shared, supabase, function-middleware, policies, otp, ...
  features/{auth,accounts,team-accounts,admin,notifications}
  billing/{core,gateway,stripe}
  cms/                      27 paquetes @pymekit/cms-*: api, auth, data-explorer,
                            filters, query-builder, audit-logs, dashboards, ...
                            y sus paquetes cliente *-ui
docs/tfg/                   plan, requisitos, decisiones, bitácora, trazabilidad
scripts/tfg/                controles del TFG (marca, comentarios, etiquetas)
.claude/                    harness de desarrollo asistido (skills, agentes, hooks)
memoria/                    esta memoria
\end{lstlisting}
  \caption{Estructura simplificada del repositorio.}
  \label{fig:imp-arbol}
\end{figure}

Las convenciones más importantes, recogidas en las directrices del
repositorio (\texttt{AGENTS.md}), son:

\begin{itemize}
  \item las \textit{server functions} se definen en ficheros
        \texttt{*.functions.ts} con la cadena
        \texttt{.middleware(\ldots).validator(Esquema).handler(\ldots)};
  \item los módulos solo de servidor llevan el sufijo \texttt{.server.ts} y
        un \textit{plugin} de Vite impide que lleguen al navegador;
  \item en el CMS, cada funcionalidad es un paquete con un servicio Drizzle
        (clase y fábrica \texttt{createXService()}), sus rutas Hono y
        funciones puras en \texttt{utils/};
  \item ningún texto visible se escribe en el código (todo pasa por i18n), y
        los registros usan el \textit{logger} compartido, nunca
        \texttt{console.log};
  \item los puntos clave se marcan con comentarios \texttt{[TFG] RF-xx} que
        un script recopila en la matriz de trazabilidad (275 etiquetas en su
        última generación).
\end{itemize}

% ---------------------------------------------------------------------
\section{Reutilización: qué se ha heredado y qué es propio}
\label{sec:imp-reutilizacion}

% PENDIENTE (P-06): redacción neutra; confirmar con el tutor la forma de
% presentar la base reutilizada sin nombrarla (ADR-007).
Como se decidió al planificar el proyecto (ADR-001), PymeKit no se ha escrito
desde cero. Parte de dos bases de código existentes, utilizadas conforme a su
licencia: una \textbf{base SaaS} sobre la misma pila tecnológica (aplicación
web, autenticación, cuentas y equipos, facturación y panel de
super-administración) y un \textbf{CMS de administración de datos} para el
mismo tipo de base de datos (aplicación de página única con su propia API y
su esquema de base de datos). Ambas se han \emph{integrado} en un único
producto, \emph{adaptadas} a sus requisitos y \emph{extendidas} con
funcionalidad nueva. El origen está documentado en el repositorio, que es
privado, y todo el código importado ha seguido el mismo flujo: copiar,
retirar las marcas originales, comentar en español, verificar y registrar.

% PENDIENTE (P-06)
La tabla~\ref{tab:imp-reutilizacion} distingue las tres categorías. La
aportación del TFG no está en reescribir lo que ya existía, sino en la
integración de dos sistemas con arquitecturas distintas, en la corrección de
su seguridad, en la adaptación al caso de una pyme y en la verificación y la
documentación del conjunto.

\begin{longtable}{@{}>{\raggedright\arraybackslash}p{2.0cm}>{\raggedright\arraybackslash}p{12.2cm}@{}}
  \caption{Código reutilizado, adaptado y propio.}\label{tab:imp-reutilizacion}\\
  \toprule
  \textbf{Tipo} & \textbf{Elementos} \\
  \midrule
  \endfirsthead
  \toprule
  \textbf{Tipo} & \textbf{Elementos} \\
  \midrule
  \endhead
  \bottomrule
  \endfoot
  Reutilizado & Biblioteca de componentes de interfaz; clientes de base de datos; \textit{middleware} de las \textit{server functions}; módulos de autenticación, cuentas, equipos, notificaciones, correo y facturación; esquema \texttt{public} de cuentas, roles, invitaciones y facturación con sus pruebas pgTAP; servicios y rutas de la API del CMS; utilidades del CMS (filtros, constructor de consultas, formateadores). \\
  \midrule
  Adaptado & Aplicación web (sin el CMS de contenidos basado en ficheros, el registro de cambios ni la documentación de la base; con comprobación de sesión viva); esquema del CMS renombrado a \texttt{cms} y endurecido (ADR-012, ADR-015); migración de unos 64 formularios y unos 183 ficheros con textos del CMS a la pila de la web (ADR-013); reescritura de las aproximadamente treinta pantallas del CMS como rutas de la web; pruebas E2E heredadas; internacionalización (solo español, ADR-021). \\
  \midrule
  Propio & Integración del CMS en la web (montaje de la API, llamada en proceso durante el SSR, guardas de la consola; ADR-011, ADR-016); pegamento entre super-administrador y rol raíz del CMS (ADR-014); corrección de 38 fallos heredados, 25 de ellos de seguridad (bitácora); blog en la base de datos y rutas públicas (ADR-017); instantánea del autor en la auditoría (ADR-018); cuadrícula propia de los paneles (ADR-019); consola por áreas de negocio, nombres en español como datos, ayuda contextual «?» y resúmenes (ADR-020); bloque académico configurable; identidad visual; 25 de las 26 migraciones añadidas durante el proyecto (la primera incorpora el esquema del CMS de partida) y las pruebas pgTAP de aislamiento escritas en las revisiones de RLS; el \textit{harness} de desarrollo y los controles del TFG. \\
\end{longtable}
\todo{Confirmar con el tutor (P-06) la forma de presentar la parte
reutilizada y, si procede, cuantificarla (ficheros o líneas reutilizados,
adaptados y nuevos).}

% ---------------------------------------------------------------------
\section{Construcción por incrementos}
\label{sec:imp-fases}

El sistema se ha construido en incrementos, cada uno funcional y verificado
antes de pasar al siguiente, con su rama de trabajo y su
\textit{pull request}. La tabla~\ref{tab:imp-fases} resume el contenido de
cada incremento y el estado de las pruebas a su cierre; la planificación
temporal se presenta en el capítulo de Planificación.

\begin{longtable}{@{}l>{\raggedright\arraybackslash}p{9.6cm}>{\raggedright\arraybackslash}p{3.0cm}@{}}
  \caption{Incrementos de construcción y estado de las pruebas al cerrarlos.}\label{tab:imp-fases}\\
  \toprule
  \textbf{Fase} & \textbf{Contenido} & \textbf{Pruebas al cierre} \\
  \midrule
  \endfirsthead
  \toprule
  \textbf{Fase} & \textbf{Contenido} & \textbf{Pruebas al cierre} \\
  \midrule
  \endhead
  \bottomrule
  \endfoot
  F0 & \textit{Harness} de desarrollo: directrices, plan, requisitos, guías, \textit{skills}, agentes, \textit{hooks}, controles del TFG e integración continua mínima & --- \\
  F1 & Base SaaS importada: renombrado del ámbito de paquetes a \texttt{@pymekit/*}, retirada de módulos no necesarios, base de datos local, puerto propio, CI completa & 61 unit.; pgTAP 27/396; E2E 43 \\
  F2.1 & Base de datos del CMS: esquema \texttt{cms}, pegamento super-administrador $\rightarrow$ Root, MFA obligatorio; primera revisión adversarial de RLS y corrección de dos brechas heredadas & pgTAP 57/1.550 \\
  F2.2 & API del CMS: 19 paquetes de servidor, Hono en \texttt{/api/cms}, autenticación con la sesión de la web & 1.130 unit.; 11 E2E de la API \\
  F2.3 & Base de la interfaz: consola compartida, guardas con \texttt{aal2}, cliente RPC, \textit{fetch} isomorfo & E2E 60 \\
  F2.4 & Explorador de datos: listado con filtros, orden, búsqueda, columnas y vistas; ficha con relaciones; creación, edición y borrado & 1.313 unit.; E2E 86 \\
  F2.5 & Exploradores de usuarios y de almacenamiento & E2E 102 \\
  F2.6 & Auditoría y búsqueda global; auditoría no falsificable y redactada según permisos & pgTAP 60/1.652; E2E 108 \\
  F2.6b & Blog en la base de datos gestionado desde el CMS y rutas públicas; formatos legibles & pgTAP 62/1.730; E2E 129 \\
  F2.7 & Ajustes: general, autenticación, miembros, RBAC (roles, grupos y permisos), recursos y diseñador de fichas & pgTAP 67/2.049; E2E 141 \\
  F2.8 & Paneles con \textit{widgets} de métrica, gráfica y tabla & pgTAP 2.060 \\
  F2.9 & Cierre de la integración del CMS: suite E2E completa, \textit{skills} actualizadas & E2E 149 \\
  F3a--b & Desmarcado completo, español por defecto (unas 1.770 claves), correos y plantillas en español, identidad visual, etiquetas legibles de enumerados, sesión caducada & E2E 150 \\
  F3c & Consola por áreas de negocio, nombres en español de las tablas, retirada del esquema de demostración (ADR-020); solo español (ADR-021); ayuda «?» y resúmenes; bloque académico en la página de inicio & 1.691 unit.; pgTAP 70/2.073; E2E 156 \\
\end{longtable}

Las cifras de pgTAP se expresan como ficheros/pruebas. Las fases F4
(comentarios didácticos), F5 (adaptación a pymes), F6 (validación) y F7
(despliegue) quedan fuera de este capítulo en la fecha de redacción.
\todo{Actualizar la tabla~\ref{tab:imp-fases} al cerrar F4--F7.}

% ---------------------------------------------------------------------
\section{Extractos de código significativos}
\label{sec:imp-extractos}

Se presentan cinco extractos breves que ilustran las decisiones del capítulo
de Diseño. Se han recortado los comentarios y las líneas no esenciales.

\subsection{Una política RLS condicionada al permiso}

El extracto~\ref{lst:imp-rls} muestra la política que permite borrar una
invitación. No basta con pertenecer a la cuenta: hace falta además el permiso
\texttt{invites.manage} en ella. La llamada a \texttt{auth.uid()} va dentro de
\texttt{(select \ldots)} para que PostgreSQL la evalúe una sola vez por
consulta y no por cada fila.

\begin{lstlisting}[style=pymekit, language=SQL, caption={Política de borrado de invitaciones (\texttt{07-invitations.sql}).}, label={lst:imp-rls}]
create policy invitations_delete on public.invitations
  for delete to authenticated
  using (
    has_role_on_account (account_id)
    and public.has_permission (
      (select auth.uid ()),
      account_id,
      'invites.manage'::public.app_permissions
    )
  );
\end{lstlisting}

El blog (extracto~\ref{lst:imp-blog}) combina tres mecanismos: se revocan
todos los privilegios, se concede la lectura \emph{por columnas} (la web no
necesita saber qué usuario creó una entrada) y la política solo devuelve lo
publicado. La escritura no se concede a ningún rol de la API: se hace desde
el CMS, con sus funciones que comprueban el permiso.

\begin{lstlisting}[style=pymekit, language=SQL, caption={Lectura pública del blog (\texttt{19-blog.sql}, recortado).}, label={lst:imp-blog}]
alter table public.blog_posts enable row level security;
revoke all on public.blog_posts from anon, authenticated, service_role;
grant select, insert, update, delete on table public.blog_posts to service_role;

grant select (id, slug, title, excerpt, content, status,
              published_at, category_id /* ... */)
  on table public.blog_posts to anon, authenticated;

create policy blog_posts_read_published on public.blog_posts
  for select to anon, authenticated
  using (status = 'published' and published_at <= now());
\end{lstlisting}

\subsection{Una \textit{server function} protegida por \textit{middleware}}

El extracto~\ref{lst:imp-sfn} es la función que crea una sesión de pago para
un equipo; es la que materializa la segunda barrera del caso CU-04. La cadena
de \textit{middleware} gestiona los errores, exige sesión y comprueba el
permiso \texttt{billing.manage} sobre la cuenta indicada (consultando
\texttt{has\_permission} en la base de datos) antes de ejecutar el
manejador. Los datos de entrada se validan con un esquema Zod.

\begin{lstlisting}[style=pymekit, caption={Server function de pago de equipo (\texttt{billing.functions.ts}, recortado).}, label={lst:imp-sfn}]
const billingManageMiddleware = [
  errorMiddleware,
  authMiddleware,
  withFeaturePermission('billing.manage'),
];

export const createTeamAccountCheckoutSession = createServerFn({
  method: 'POST',
})
  .middleware(billingManageMiddleware)
  .validator(TeamCheckoutSchema)
  .handler(async ({ data }) => {
    if (!featureFlagsConfig.enableTeamAccountBilling) {
      throw new Error('Team account billing is not enabled');
    }

    const client = getSupabaseServerClient(); // cliente con RLS
    const service = createTeamBillingService(client);

    return service.createCheckout(data);
  });
\end{lstlisting}

\subsection{Montaje de la API del CMS y una ruta Hono}

La integración del CMS en la web (ADR-011) se reduce a una ruta de servidor
que reenvía cualquier petición bajo \texttt{/api/cms/} a la aplicación Hono
(extracto~\ref{lst:imp-montaje}). La importación dinámica dentro del
manejador garantiza que el código del CMS ---con el controlador de PostgreSQL
y la clave secreta--- nunca forme parte de un fragmento del navegador.

\begin{lstlisting}[style=pymekit, caption={Montaje de la API del CMS (\texttt{routes/api/cms/\$.ts}).}, label={lst:imp-montaje}]
async function handleCmsApiRequest({ request }: { request: Request }) {
  const { getCmsApiApp } = await import('#/lib/cms/cms-api-app.server.ts');
  const app = await getCmsApiApp();

  return app.fetch(request);
}

export const Route = createFileRoute('/api/cms/$')({
  server: {
    handlers: {
      GET: handleCmsApiRequest,
      POST: handleCmsApiRequest,
      // ... PUT, PATCH, DELETE y OPTIONS igual
    },
  },
});
\end{lstlisting}

Cada funcionalidad del CMS registra sus rutas en la aplicación Hono
(extracto~\ref{lst:imp-hono}): valida la entrada, delega la autorización y la
consulta en su servicio y traduce cualquier error a un código estable, sin
exponer el mensaje interno.

\begin{lstlisting}[style=pymekit, caption={Ruta del registro de auditoría (\texttt{get-audit-logs-route.ts}, recortado).}, label={lst:imp-hono}]
export function registerAuditLogsRoute(router: Hono) {
  return router.get(
    '/v1/audit-logs',
    zValidator('query', AuditLogsQuerySchema, invalidAuditLogsInput),
    async (c) => {
      const { cursor, limit, ...filters } = c.req.valid('query');

      try {
        const data = await createAuditLogsService(c).getAuditLogs({
          cursor, limit, filters,
        });
        return c.json(data);
      } catch (error) {
        // 4xx/5xx con código estable (AUDIT_LOG_*); el detalle va al log
        return respondWithAuditLogsError(c, error, { route: 'list' });
      }
    },
  );
}
\end{lstlisting}

El servicio ejecuta cada consulta dentro de una transacción en la que antes
fija la identidad del usuario (extracto~\ref{lst:imp-claims}). Los valores se
pasan como \emph{parámetros enlazados}, nunca concatenados en el texto SQL, y
el rol se limita a los que puede llevar un JWT de usuario. Así, las políticas
RLS del esquema \texttt{cms} se aplican a la API igual que a cualquier otro
cliente.

\begin{lstlisting}[style=pymekit, caption={Identidad del usuario en cada transacción del CMS (\texttt{drizzle-client.ts}, recortado).}, label={lst:imp-claims}]
const role = ALLOWED_DB_ROLES.has(token.role) ? token.role : 'anon';
const claims = JSON.stringify(token);

return rlsClient.transaction(async (tx) => {
  await tx.execute(sql`select set_config('request.jwt.claims', ${claims}, true)`);
  await tx.execute(sql`select set_config('request.jwt.claim.sub', ${sub}, true)`);
  await tx.execute(sql`select set_config('role', ${role}, true)`);

  return await transaction(tx);
});
\end{lstlisting}

\subsection{La migración de las áreas de la consola}

La consola por áreas (ADR-020) no necesitó cambios de estructura: el área de
cada tabla se guarda en el campo \texttt{ui\_config} de sus metadatos, que ya
existía. La migración (extracto~\ref{lst:imp-areas}) carga los valores por
defecto en una tabla temporal y los aplica \emph{solo} donde no hay un valor
propio, de modo que nunca pisa lo que una pyme haya personalizado desde
Ajustes $\rightarrow$ Recursos y puede ejecutarse de nuevo sin efectos.

\begin{lstlisting}[style=pymekit, language=SQL, caption={Áreas de navegación del CMS (\texttt{20261002140000\_cms\_navigation\_groups.sql}, recortado).}, label={lst:imp-areas}]
insert into cms_navigation_defaults
    (schema_name, table_name, navigation_group, display_name, previous_names, ordering)
values
    ('public', 'blog_posts', 'Blog', 'Entradas', array['Blog: entradas'], 100),
    ('public', 'accounts', 'Cuentas', 'Cuentas', array['Cuentas'], 200),
    ('public', 'subscriptions', 'Facturación', 'Suscripciones',
     array['Facturación: suscripciones'], 310),
    ('public', 'config', 'Sistema', 'Configuración',
     array['Configuración de la plataforma'], 400);
    -- ... resto de tablas

-- Área: solo si la tabla todavía no tiene una.
update cms.table_metadata as meta
set ui_config = coalesce(meta.ui_config, '{}'::jsonb)
        || jsonb_build_object('navigation_group', defaults.navigation_group)
from cms_navigation_defaults as defaults
where meta.schema_name = defaults.schema_name
  and meta.table_name = defaults.table_name
  and nullif(btrim(meta.ui_config ->> 'navigation_group'), '') is null;
\end{lstlisting}

% ---------------------------------------------------------------------
\section{Desarrollo asistido por agentes de IA}
\label{sec:imp-agentes}

El proyecto se ha desarrollado con la asistencia de un agente de
programación de IA (Claude Code~\cite{claudecode}). El autor fija el alcance,
toma las decisiones y revisa los resultados; el agente implementa, prueba y
documenta bajo unas reglas explícitas. Para que esa asistencia sea
controlable y reproducible se construyó, antes de escribir código, un
\textbf{\textit{harness}}: el conjunto de directrices, procedimientos y
controles automáticos que gobiernan el trabajo del agente
(tabla~\ref{tab:imp-harness}).

\begin{table}[htbp]
  \centering
  \footnotesize
  \caption{Componentes del \textit{harness} de desarrollo asistido.}
  \label{tab:imp-harness}
  \begin{tabularx}{\textwidth}{@{}>{\raggedright\arraybackslash}p{3.0cm}X@{}}
    \toprule
    \textbf{Componente} & \textbf{Contenido} \\
    \midrule
    Directrices & \texttt{AGENTS.md} raíz (fuente única) y uno por aplicación o paquete: pila, patrones, reglas de seguridad de la base de datos, convenciones e idioma, y una lista de verificación obligatoria al terminar cada cambio \\
    Registros & Plan, progreso, requisitos, decisiones (ADR), bitácora de incidencias, trazabilidad y mapa de reutilización en \texttt{docs/tfg/}, que el agente actualiza en cada cambio \\
    \textit{Skills} & Procedimientos reutilizables: incorporar un módulo de una base de partida, comentar un módulo en español, registrar una decisión, redactar una sección de la memoria, revisión adversarial del \textit{diff}, búsqueda de fallos, revisión de RLS, experto en PostgreSQL y patrones de servicios, formularios y E2E \\
    Subagentes & Revisor de calidad, revisor de comentarios y redactor de la memoria, con instrucciones y herramientas propias \\
    \textit{Hooks} & Al iniciar la sesión, muestran la fase actual y sus pendientes; tras cada edición, ejecutan el control de marca sobre el fichero; antes de terminar, exigen actualizar el progreso si hay cambios de código \\
    Permisos & Lectura permitida y escritura denegada sobre las bases de partida; lectura denegada sobre ficheros de secretos \\
    Controles en CI & Control de marca bloqueante e informe de comentarios en inglés; un script aparte regenera la tabla de etiquetas de trazabilidad \\
    \bottomrule
  \end{tabularx}
\end{table}

El procedimiento de revisión de RLS merece mención aparte porque es el que
más hallazgos ha producido. Su premisa es que leer las políticas no basta:
hay que \emph{demostrar} el aislamiento. Tras un inventario y una auditoría
estática con un catálogo de ataques, cada hallazgo plausible se convierte en
una prueba pgTAP que ejecuta el ataque como el usuario atacante y comprueba
que falla; después, un agente distinto, que solo recibe la política, la
prueba y la afirmación, intenta \emph{refutarla}. Una brecha solo se acepta
si el refutador la reproduce, y un aprobado, solo si no encuentra un vector
abierto. La refutación encontró, por ejemplo, la lectura del catálogo de
roles de PostgreSQL que la primera auditoría había dado por buena (B-08). Sus
resultados se presentan en el capítulo de Pruebas.

La experiencia también mostró los límites del enfoque. La verificación
exhaustiva con varios agentes por paso (implementador, revisor de seguridad y
refutador, y la suite E2E completa en cada paso) encontró brechas reales
(B-45, B-47), pero su coste en \textit{tokens} resultó excesivo, y el autor
pidió reducirlo (B-48). Desde entonces se trabaja en un «modo ahorro»: un
agente por paso, E2E completos solo al cerrar cada fase y revisión de RLS
solo en cambios importantes de permisos. El riesgo se traslada al cierre de
fase, como mostraron dos pruebas desactualizadas que solo detectó la suite
completa (B-51, B-54). Además, ninguna revisión automática es exhaustiva: una
inyección SQL heredada pasó por varias revisiones de agentes sin ser vista y
la encontró la revisión crítica de sus informes (B-24), y algunos fallos solo
los descubrió el autor con pruebas manuales (B-14, B-40, B-52).
\todo{Valorar en las Conclusiones el efecto de la asistencia de IA con los
datos de dedicación (horas del autor frente a tiempo de reloj).}

% ---------------------------------------------------------------------
\section{Problemas relevantes resueltos}
\label{sec:imp-problemas}

Todas las incidencias se registraron en una bitácora en el momento de
producirse, con el formato «qué pasó, causa, solución, evidencia y lección».
Hasta la fase F3c se registraron 55. La tabla~\ref{tab:imp-bitacora} recoge
una selección de las no relacionadas directamente con el control de acceso;
las de seguridad se analizan en el capítulo de Pruebas
(sección~\ref{sec:pru-seguridad}).

\begin{longtable}{@{}l>{\raggedright\arraybackslash}p{4.8cm}>{\raggedright\arraybackslash}p{7.4cm}@{}}
  \caption{Selección de problemas resueltos y lecciones aprendidas (bitácora del proyecto).}\label{tab:imp-bitacora}\\
  \toprule
  \textbf{ID} & \textbf{Problema} & \textbf{Solución y lección} \\
  \midrule
  \endfirsthead
  \toprule
  \textbf{ID} & \textbf{Problema} & \textbf{Solución y lección} \\
  \midrule
  \endhead
  \bottomrule
  \endfoot
  B-01 & El renombrado masivo del ámbito de paquetes no alcanzó una expresión regular donde aparecía escapado & Se corrigió y el control de marca detecta ahora también las variantes escapadas. Un renombrado masivo se verifica buscando variantes, no solo la forma literal. \\
  B-04 & La base de datos local quedó accesible desde Internet porque los contenedores se saltan el cortafuegos & Regla de filtrado específica para contenedores y acceso por túnel SSH (ADR-010). En un servidor con IP pública, «local» no significa privado. \\
  B-10 & Esquema declarativo y migraciones del CMS de partida desincronizados (nueve funciones distintas) & Se alinearon con las migraciones. La comparación de esquemas forma parte de la verificación de cualquier cambio de base de datos. \\
  B-11 & Una prueba heredada se titulaba «MFA obligatorio por defecto», pero su aserción comprobaba que se concedía el acceso & Se reescribió. Una prueba en verde solo demuestra que el sistema hace lo que dice la prueba: hay que leer las aserciones. \\
  B-16 & Ciclo de dependencias al poner componentes y API del explorador en el mismo paquete & Paquete cliente separado. Separar cliente y servidor también evita ciclos. \\
  B-17, B-29 & La API respondía 500 a registros inexistentes porque el envoltorio de errores perdía el código de PostgreSQL & Se conserva la causa y se traduce el código SQLSTATE a 403/404. Sin la causa no se distingue «no existe» de «ha fallado». \\
  B-35 & La herramienta de comparación de esquemas omitía la opción \texttt{security\_invoker} de una vista y los permisos por columna & Esa parte de la migración se escribe a mano y pgTAP comprueba ambas propiedades. Que la comparación salga limpia no demuestra que la migración sea segura. \\
  B-36 & La sincronización de tablas del CMS duplicaba las relaciones en cada ejecución & Función idempotente, cubierta por pgTAP. \\
  B-40, B-46 & El control de marca solo miraba el contenido de los ficheros, no sus rutas, y no podía ver la marca dibujada en imágenes & Se comprueban también las rutas; logotipo y capturas propias generadas con un script. Un control automático solo cubre lo que se le dice que mire. \\
  B-52 & Tras cambiar de rama, el servidor de desarrollo conservaba un árbol de rutas sin el CMS & Reinicio limpiando cachés. Que las pruebas pasen contra la \textit{build} no garantiza que el entorno de desarrollo esté sano. \\
  B-53 & Una sesión revocada pero con firma válida rompía la consola & Comprobación con estado de la sesión en la entrada de las rutas protegidas. Verificar la firma de un JWT no demuestra que la sesión exista. \\
  B-55 & La lista de miembros superaba el tiempo máximo con unos cientos de cuentas (función volátil evaluada por fila) & Envolver la llamada en \texttt{(select \ldots)}; véase la sección~\ref{sec:pru-rendimiento}. \\
\end{longtable}
